
\section{Bessel Process, simulation}
\subsection{Numerical Scheme for SDE}

The simplest numerical scheme for simulating SDEs is the Euler-Maruyama method, described in \autocite[Chp.~9,10]{kloedenNumericalSolution1992}.
It's higher-order correction is the Milstein scheme.

The Euler-Maruyama method is a simple extension of the Euler method for ODEs.
Given an SDE
\begin{equation}
    dX_t = b(X_t, t) dt + \sigma(X_t, t) dW_t,
\end{equation}
the Euler-Maruyama method is
\begin{equation}
    X_{n+1} = X_n + b(X_n, t_n) \Delta t + \sigma(X_n, t_n) \Delta W_n,
\end{equation}
where $\Delta t = t_{n+1} - t_n$ and $\Delta W_n = W_{n+1} - W_n$.

The Milstein scheme is a higher-order correction to the Euler-Maruyama method.
It is given by
\begin{equation}
    X_{n+1} = X_n + b(X_n, t_n) \Delta t + \sigma(X_n, t_n) \Delta W_n + \frac{1}{2} \sigma(X_n, t_n) \sigma'(X_n, t_n) (\Delta W_n^2 - \Delta t),
\end{equation}
where $\sigma'(X_n, t_n)$ is the derivative of $\sigma$ with respect to $X$ evaluated at $X_n$ and $t_n$.

The Euler-Maruyama method is essentially replacing the integral with a Riemann sum, using the left endpoint of each interval. It converges to the Ito integral of the SDE.

\ask{What is the difference for the Stratonovich integral?}

\subsection{Bessel Process}

Example 4.2.2 Let $B=\left(B_1, \ldots, B_n\right)$ be Brownian motion in $\mathbf{R}^n, n \geq 2$, and consider

\[
R(t, \omega)=|B(t, \omega)|=\left(B_1^2(t, \omega)+\cdots+B_n^2(t, \omega)\right)^{\frac{1}{2}}
\]

i.e. the distance to the origin of $B(t, \omega)$. The function $g(t, x)=|x|$ is not $C^2$ at the origin, but since $B_t$ never hits the origin, a.s. when $n \geq 2$ (see Oksendal Exercise 9.7, or SDE lecture note) Itô's formula still works and we get

\[
d R= \frac{B(t)}{R} \cdot d B(t)+\frac{n-1}{2 R} d t
\]

\begin{intuition}
    The first term: $B(t) / R$ is the unit vector pointing to current direction; when $R$ is large, it changes according to the Brownian motion's $R$-direction.

    The second term: When $R$ is small, it doesn't want to reach the origin, so we have an effective drift away from the origin. \todo{The exact reason why we have drift to the order of $1/R$ is not clear to me.}
\end{intuition}